\chapter{机器学习究竟在学习什么}

\begin{quote}\small
\textit{本章摘要。}本章定义样本、任务、模型、损失和风险，说明不同损失对应不同预测对象，并讨论正则化、可识别性及预测到决策的边界。
\end{quote}

同一段振动记录，可以用来判断轴承是否故障，也可以用来预测下一时刻的振幅。这两个问题需要不同的标签，预测错误的后果也不同。上一章介绍了计算所需的数学工具；本章先确定要学习的对象，再说明怎样用目标函数比较模型。

输入决定模型能够使用哪些信息，标签规定需要预测什么，损失衡量预测错误的代价。如果三者与实际任务不一致，即使训练目标已经达到最小值，预测结果也可能无法支持维护决策。本章由数据和任务引出模型、损失与风险，再讨论归纳偏置、正则化和任务的可学习性。后续各章将把这些概念用于具体计算。

\section{学习对象：数据、任务与模型}

% auxiliary-resource-guide
本章末的“辅助学习材料”列出与核心方法对应的讲解和实践入口，可在阅读相关推导后，按所列小节对照学习。

\paragraph{从数据到任务。}
数据（data）是对对象、过程或环境的观测记录。一个样本（sample）是一次相对完整的观测，通常记为$\bm{x}\in\R^d$；监督学习中还伴随标签$y$，于是样本写为$(\bm{x},y)$。这里的“样本”不是数据表的一行这么简单：一次设备运行周期、一张图像、一个句子或一段传感器窗口都可以成为样本，关键在于它是否对应一个明确的可学习单位。

“任务”由输入空间$\mathcal{X}$、输出空间$\mathcal{Y}$以及评价准则共同决定。输入空间是模型允许接收的对象集合（例如时序数据、图像、文本等），输出空间是允许回答的结果集合，这两个集合先于可学习参数的确定。分类的输出是离散类别，回归的输出是连续数值，预测性维护可能输出剩余寿命，异常检测则可能没有逐样本标签。相同的数据可以支持不同任务。例如，振动信号既可以用于故障分类，也可以用于生成未来波形或估计设备状态。

\paragraph{把一个任务写成可核查的样本。}
假设每小时读一次温度和振动，把这两项读数排成输入向量：
\[
 \bm x=(\text{温度},\text{振动})^{\mathsf T}\in\R^2.
\]
标签为下一小时温度$y\in\R$时，这就是回归任务；改成“未来一天是否故障”的$y\in\{0,1\}$时，就成为二分类。

监督学习用已知输入和标签学习对应关系，无监督学习则在没有这些目标标签时寻找输入结构，即发现样本之间的相似性、变量之间的关联等规律。例如，不提供故障标签，仅根据温度和振动把相似的设备记录分组，但这些分组不自动代表正常或故障。

同一数据记录能面向不同任务，不同任务同样需要不同的标签或指标。训练是利用已有样本调整参数；推理是在参数训练固定后处理一个新输入，得到对新样本标签或指标的预测。


\paragraph{模型、参数与预测。}
模型规定了允许采用的预测关系。直接给出一个预测值时，可以用确定性模型（deterministic model）表示：
\begin{equation}
 f_{\bm{\theta}}:\mathcal{X}\rightarrow\mathcal{Y},
 \qquad \bm{\theta}\in\Theta.
\end{equation}
这里$\Theta$是允许的参数取值集合。给定输入$\bm{x}$和参数$\bm\theta$，输出为$\hat y=f_{\bm\theta}(\bm x)$。模型结构规定函数的形式，参数从这些函数中确定一个具体函数。例如，$f_{w,b}(x)=wx+b$规定了直线形式，$w$和$b$分别决定斜率和截距。训练利用数据和目标函数选择参数，求解算法决定怎样完成这一选择。

如果还需要描述不同标签的可能性，可以建立概率模型（probabilistic model），输出条件分布$p_{\bm\theta}(y\mid\bm x)$。在$K$类分类中，输出是$K$个非负且总和为一的概率。在高斯回归中，模型可以给出条件均值和方差。概率分布随后可以转成一个预测值，例如取回归分布的均值，或取概率最大的类别；这种转换会舍去部分不确定性信息。确定性预测也可以表示为集中在单点$f_{\bm\theta}(\bm x)$上的退化分布，但它没有表达该点以外的观测波动。

上述两类表述都直接描述给定输入后的输出，属于判别建模（discriminative modeling）。选择函数形式时，需要考虑输入结构、任务和计算资源。例如，处理图像时可以利用邻近像素的局部关系，处理时间序列时可以利用观测的先后顺序。具体网络结构将在相应章节介绍。

概率建模还可以作用于未观测的状态或未知参数。用潜变量$z$表示设备状态时，可通过$p_{\bm\theta}(z\mid\bm x)$和$p_{\bm\theta}(y\mid\bm x,z)$描述预测，再对$z$求和或积分。贝叶斯模型则用后验分布$p(\bm\theta\mid\mathcal D)$表达观测训练数据后对参数的不确定性。因此，输出是否为概率分布，与参数是否用分布表示，是两个需要分别确定的问题。

生成模型（generative model）描述观测数据的分布。在带标签任务中，它可以建立联合分布$p_{\bm\theta}(\bm x,y)$；没有标签时，也可以直接建立$p_{\bm\theta}(\bm x)$。联合模型能够通过条件概率得到$p_{\bm\theta}(y\mid\bm x)$，也可能用于生成新样本。相比直接建立条件分布，它增加了关于输入如何产生的假设，这些假设同样需要检验。

\section{损失函数怎样规定训练目标}

\paragraph{从一次预测的误差到整个训练集的目标。}
损失函数（loss function）把模型输出与观测标签之间的差异转成一个数。数值越大，表示该预测按所选规则付出的代价越高。例如，回归任务可以使用平方损失：
\begin{equation}
 \ell_{\mathrm{sq}}\bigl(f_{\bm{\theta}}(\bm{x}),y\bigr)=\frac12\bigl(f_{\bm{\theta}}(\bm{x})-y\bigr)^2.
\end{equation}
这里$f_{\bm\theta}(\bm x)$是预测值，$y$是观测标签，系数$1/2$使求导更简洁。平方损失会放大较大的误差。如果任务对高估和低估的代价不同，还需要改变损失形式，第四章将讨论这种选择。

对训练集$\mathcal D=\{(\bm x_i,y_i)\}_{i=1}^n$，单样本损失的平均值称为经验风险（empirical risk）：
\begin{equation}
 \hat R(f_{\bm\theta})=\frac1n\sum_{i=1}^n\ell(f_{\bm\theta}(\bm x_i),y_i).
\end{equation}
训练目标函数（objective function）规定需要最小化或最大化的量。它可以是经验风险，也可以包含正则项。单样本损失、经验风险和完整训练目标描述不同层次的量，需要分别说明。

概率模型可以通过似然构造目标。把参数视为固定未知量时，极大似然估计（Maximum Likelihood Estimation，ML）选择使观测数据似然最大的参数。在给定参数后样本相互独立的联合模型中，有
\begin{equation}
 \hat{\bm{\theta}}_{\mathrm{ML}}
 =\argmax_{\bm{\theta}}\log p_{\bm{\theta}}(\mathcal D)
 =\argmax_{\bm{\theta}}\sum_{i=1}^{n}\log p_{\bm{\theta}}(\bm{x}_i,y_i).
\end{equation}
条件模型则使用$\sum_i\log p_{\bm\theta}(y_i\mid\bm x_i)$。最大化对数似然等价于最小化其负值，称为负对数似然（Negative Log-Likelihood，NLL）。第三章将说明高斯回归的NLL为何对应平方误差。

如果为参数指定先验$p(\bm\theta)$，还可以选择后验密度最大的参数，称为最大后验估计（Maximum A Posteriori estimation，MAP）：
\begin{equation}
 \hat{\bm{\theta}}_{\mathrm{MAP}}
 =\argmax_{\bm{\theta}}p(\bm{\theta}\mid\mathcal D)
 =\argmax_{\bm{\theta}}\bigl[\log p(\mathcal D\mid\bm{\theta})+\log p(\bm{\theta})\bigr].
\end{equation}
该等式要求后验可归一化；其归一化常数与待比较的参数无关。MAP使用先验，但结果仍是一个参数点。保留完整后验分布还需要进一步计算。极大似然和MAP都可以用于条件模型或联合模型，采用哪一种估计方法与是否建立输入的分布是两个不同的选择。

\paragraph{经验风险与总体风险.}  
有限样本上的拟合与总体上的预测能力需要分开讨论，经验分布到真实分布的一致收敛是统计学习理论的基础问题之一\cite{vapnik1971uniform}。
设$(X,Y)\sim P$描述部署环境中的输入与标签，$X\in\R^d$，$Y$可为实数或类别。总体风险定义为
\begin{equation}
 R(f)=\E_{(X,Y)\sim P}[\ell(f(X),Y)].
\end{equation}
这里的期望包含尚未观察到的未来样本，因此通常不能直接计算。若训练样本独立同分布于$P$（independent and identically distributed，彼此独立且共享同一个分布），对一个事先固定且损失可积的$f$，有$\E[\hat R(f)]=R(f)$。训练前允许选择的模型参数属于参数空间$\Theta$；“事先固定”强调不能先看测试结果再挑选$f$而仍使用这条无偏性解释。若损失方差有限，独立性进一步给出
\begin{equation}
 \Var(\hat R(f))=\frac{1}{n^2}\sum_{i=1}^n\Var(\ell(f(X_i),Y_i))
 =\frac{\Var(\ell(f(X),Y))}{n}.
\end{equation}
这解释了平均损失为什么可用于估计风险，也暴露了限制：相邻设备窗口相关时有额外协方差项；$f$由同一训练集选择后，上述固定模型的无偏论证不能直接用于其训练误差。验证与测试的必要性由此产生，而不是来自惯例。

训练时通常最小化经验风险，真正希望最小化的却是部署分布上的总体风险。“\textbf{过拟合}”的核心表现，是模型在训练集上的经验风险很小，但在未参与训练的新样本上的总体风险仍然较大，因而出现较大的泛化差距。模型容量过大时，它不仅可以学习输入与标签之间可迁移的规律，也可能记住训练样本中的噪声、偶然相关性甚至标签错误；但模型容量过大并不是过拟合的必要条件，训练样本过少、标签噪声较强或反复利用验证集也可能造成过拟合。当训练集有限或含有噪声时，模型记住的细节可能无法在同一分布的新样本中重复出现。若训练集与部署分布不同，则还存在分布偏移问题，不能仅把性能下降归因于过拟合。典型地，训练早期经验风险和总体风险可能同时下降；训练后期经验风险继续下降，而总体风险可能转而上升，因为模型开始更多地拟合训练数据中的偶然性。


\section{归纳偏置}

前文提到了过拟合，通俗来讲指的是训练误差很低而新样本误差较高。与之对应的一个概念叫做“\textbf{泛化性}”，指模型在“未参与训练和选择的数据”上的有效程度。使用损失函数或目标函数进行训练，直接下降的是经验风险；但我们真正关心的是部署分布上的总体风险。经验风险最小化并不自动意味着总体风险最小化，因此需要在模型构造和训练目标中加入能够约束学习过程的偏好。下面从模型角度和训练目标角度分别说明这种约束如何发挥作用。

模型结构、参数化方式和训练算法会限制能够采用的预测关系，也会使某些解更容易被找到。这些预先引入的限制和偏好称为归纳偏置（inductive bias）。当训练样本不足以唯一确定预测关系时，归纳偏置提供选择候选解的依据。它能否改善新样本上的表现，取决于这些假设是否适合任务。

例如，对同一图像分类任务，卷积神经网络把局部连接和参数共享写进了模型结构，因此偏好局部模式以及对平移具有一定稳定性的函数；全连接网络则没有这种结构性限制。若图像任务确实具有这些性质，卷积网络就能把有限样本用于学习更有用的变化，而不必分别记住每个像素位置的完全独立关系，因此通常更容易获得较好的泛化性能。这里的“更适合”是统计意义上的，不是保证在任何数据集上都更好。

归纳偏置与可解释性有关，但二者不是同一个概念。可解释性讨论人能否理解模型的预测依据；归纳偏置讨论模型在学习时预先接受了哪些结构假设。一个模型可以具有明确的归纳偏置，却仍然难以解释；也可以相对容易解释，但其归纳偏置并不适合当前任务。判断归纳偏置是否良好，关键在于它是否与数据生成过程和任务目标相匹配，并能否在有限样本下缩小训练风险与总体风险之间的差距。

正则化（regularization）提供了加入这种偏好的另一条途径：不一定更换模型结构，而是在训练目标中明确惩罚某些不希望出现的解。其基本形式是
\begin{equation}
 \min_{\bm{\theta}}\;\hat{R}(\bm{\theta})+\lambda\Omega(\bm{\theta}),
 \qquad \Omega(\bm{\theta})=\frac12\|\bm{\theta}\|_2^2.\label{eq:foundations-regularized-risk}
\end{equation}
这里第一项$\hat R(\bm{\theta})$要求模型解释训练数据，第二项$\Omega(\bm{\theta})$表达我们对模型复杂度的偏好。选择$\Omega(\bm{\theta})=\frac12\|\bm{\theta}\|_2^2$，表示在其他条件相近时偏好范数较小的参数；它并不声称真实参数一定接近零，而是把“大权重可能对应不稳定拟合”作为一种可调的建模偏好。系数$\lambda\geq0$决定这种偏好有多强：$\lambda=0$时只拟合训练数据；$\lambda$增大时，模型愿意牺牲一部分训练集拟合精度，以换取更小的参数和可能更好的泛化。正则化的目的不是让训练误差尽可能低，而是在经验风险和对复杂度的约束之间寻找适合部署的平衡。

\begin{example}
 根据温度记录预测下一时刻温度时，若已知相邻时刻通常连续变化，可以采用平滑性假设；若设备存在周期，可以把相应周期特征加入输入。这些选择限制了候选预测关系。只有假设与设备运行条件相符时，这种限制才可能帮助有限样本下的学习。
\end{example}

\paragraph{惩罚形式与约束形式如何联系。}
为什么可以把正则化理解成“控制复杂度”？一种直接的写法是先规定允许的复杂度上限：
\begin{equation}
 \min_{\theta}\hat R(\theta)\quad\text{subject to}\quad \Omega(\theta)\leq c.
\end{equation}
这句话的意思是：在所有不超过复杂度预算$c$的模型中，选择训练风险最小的模型。它明确表达了两件事：训练数据需要拟合，但模型不能为了拟合少量样本而无限增加复杂度。

约束问题也可以通过拉格朗日乘子写成
\begin{equation}
 L(\theta,\lambda)=\hat R(\theta)+\lambda\bigl(\Omega(\theta)-c\bigr),\qquad \lambda\geq0.
\end{equation}
固定$\lambda$后，$-\lambda c$与$\theta$无关，不会改变参数最优解，于是得到惩罚形式$\hat R(\theta)+\lambda\Omega(\theta)$。因此，$c$和$\lambda$分别是两种不同的控制方式：前者规定复杂度上限，后者规定超出较简单解的代价。在线性凸问题并满足适当严格可行性条件时，二者可以通过拉格朗日最优性条件联系起来；对非凸网络，不能据此声称任意约束问题都与某个惩罚问题全局等价，也不能认为$\lambda$与$c$必然一一对应。

在约束问题的最优点$\hat\theta$处，若目标与约束满足相应的凸性条件，可以写出
\begin{equation}
 0\in\partial\hat R(\hat\theta)+\lambda\partial\Omega(\hat\theta),\qquad
 \Omega(\hat\theta)\leq c,\qquad
 \lambda(\Omega(\hat\theta)-c)=0.
\end{equation}
$\partial$是允许不可微凸函数的次梯度集合。最后一个等式是互补条件。若$\Omega(\hat\theta)<c$，说明复杂度预算没有用完，限制没有起作用，此时可以有$\lambda=0$；若$\Omega(\hat\theta)=c$，说明模型已经碰到复杂度边界，进一步降低训练风险通常需要增加复杂度，约束因此会通过正的$\lambda$影响最优解。这里的“补偿”不是说两个数必须相等，而是说训练风险的改善必须值得为增加复杂度所付出的代价。

最后用一维例子看清“收缩”究竟发生了什么。设无正则化时的训练风险为
\begin{equation}
 \hat R(w)=\frac12(w-a)^2,
\end{equation}
其中$a$是只根据训练数据得到的最优估计。它的最小值在$w=a$处，因为此时$w$与$a$完全一致。现在加入平方惩罚，训练目标变为
\begin{equation}
 J(w)=\frac12(w-a)^2+\frac{\lambda}{2}w^2.
\end{equation}
第一项希望$w$接近数据给出的估计$a$，第二项希望$w$接近零。对$w$求导得到
\begin{equation}
 J'(w)=(w-a)+\lambda w.
\end{equation}
令导数为零，得到$(1+\lambda)w=a$，因此
\begin{equation}
 \hat w=\frac{a}{1+\lambda}.
\end{equation}
由于$\lambda\geq0$，有$|\hat w|\leq|a|$；当$\lambda=0$时不收缩，随着$\lambda$增大，$\hat w$逐渐靠近零。这就是“收缩”的含义：正则化不会自动判断某个变量是噪声，也不会直接把所有不重要的系数删掉，而是把参数估计向预先选定的中心（这里是零）拉近。

收缩为什么可能改善总体风险？训练数据中的偶然波动会使无正则化估计$a$在不同训练集之间变化很大，这种变化称为估计的不稳定性或方差。平方惩罚牺牲了一部分对当前训练集的贴合，使估计更稳定；但如果真实系数确实远离零，收缩会带来偏差。因此，正则化是在偏差与方差之间作权衡：只有方差的降低足以抵消偏差增加，使用验证数据选择的$\lambda$才可能降低总体预测风险。它不是“惩罚越强越好”，也不是“参数越小越接近真实机制”。

模型选择至少同时受四类约束：数据规模与标注质量、任务中的依赖结构、训练与推理资源、错误的业务代价。复杂模型并不自动优于简单模型；它只有在表示了简单模型无法表示的结构时才有必要。后文讨论的序列、视觉和语言模型，实际上是在不同结构约束下对这四个问题的回答。


\section{任务的可学习性}
在选择模型之前，还应先问一个更根本的问题：给定的输入中，是否含有足以预测目标的稳定信息？例如，设输入是风电机组叶片的图像，标签是同一时刻的发电功率。凭借领域知识可以预期，功率还强烈依赖风速、风向、空气密度、机组运行状态和控制策略；一张只反映叶片外观的图像，通常不足以充分决定功率。图像中也许包含结冰、污损或损伤等与功率有关的间接线索，但这种联系必须由数据和机理共同证实，不能因为两个变量出现在同一张数据表中，就假定二者必然存在可用于预测的关系。

训练集上的拟合能力不能单独回答“任务是否可以被学习”这个问题。只要模型足够灵活，它就可以把输入空间划分为许多很小的区域，并为每个训练样本建立对应映射。以决策树回归为例，当训练输入彼此可区分，且树的深度、叶节点样本数等不受限制时，树可以把训练样本逐个分到不同叶节点，从而达到零训练误差；分类树也可以达到百分之百训练准确率。这说明模型有能力记住训练样本，却不说明图像中存在对未来功率稳定有效的预测线索。若新图像落入训练时未被充分覆盖的区域，模型只能依赖未经验证的划分和映射，训练集上的完美表现便可能无法迁移到新样本。这正是过拟合的一种典型表现。

从统计角度看，需要检查给定输入$X$后的条件分布$P(Y\mid X)$是否具有可利用的结构，以及这种结构在训练和部署环境之间是否足够稳定。相近输入对应相近的标签分布，是一种可能的结构；连续性、平滑性和低维结构也可以支持对未见样本的预测。但这些条件并不适用于所有任务，判断时还要结合样本覆盖、噪声水平和模型假设。即使存在稳定关系，有限训练数据也未必足以学到它。

现实中的困难在于，输入空间往往高维，且样本真实分布不均匀；有限样本更不可能把整个空间都覆盖。因此，数据中是否存在规律通常不能仅靠观察几个样本来断言。一个有用的直觉是：若相同或非常相近的输入在输入空间的一个小邻域内反复对应差异很大的标签，那么只使用这些输入时就存在不可消除的不确定性；即使知道真实条件分布，也不可能准确预测每一个样本。相反，若$X$能够稳定地缩小$Y$的不确定性，使合适的预测函数在独立新样本上显著优于“不看输入、只预测平均值或多数类”的基线，则该任务具有可学习的信号。这里的“规律”不要求是线性的、显式的，或容易写成公式；它可以是高度非线性的统计关系，但必须能够在新的数据中重复出现。

有限数据无法从逻辑上证明这种规律一定存在，只能不断积累证据。实践中应结合四类检查：首先用领域机理审查输入到标签的路径，确认不是把结果发生之后的信息误当作输入；其次检查标签的测量质量、时间对齐和样本覆盖，避免真实信号被噪声或错误配对淹没；再次在严格隔离的验证集或测试集上，与不使用输入的简单基线比较；最后检验结果是否能跨时间、机组或运行条件复现。若模型只在训练集上优于基线，或一旦按时间、设备划分测试集就失效，更合理的结论通常是现有输入不足、分布发生了变化，或模型只记住了偶然性，而不是急于增加模型复杂度。

\paragraph{从训练到部署及反馈。}
沿图\ref{fig:foundations-learning-loop}从采集到上线的路径，可以依次看四个阶段。先把设备记录整理成输入和标签，再选择能处理这些输入的模型，通过训练求出参数，最后用这些参数预测新记录。上线后的错误又会提示我们补数据、改标签或重新训练。评价时除了训练损失，还要检查新工况下的表现、概率是否校准、计算是否及时，以及误报和漏报的后果。

\begin{figure}[htbp]
\centering
\begin{tikzpicture}[>=Latex, node distance=0.85cm, every node/.style={font=\small}]
 \node[draw, rounded corners, fill=blue!10, minimum width=2.15cm, minimum height=0.9cm] (data) {\shortstack{数据\\$\mathcal D$}};
 \node[draw, rounded corners, fill=orange!12, right=of data, minimum width=2.35cm, minimum height=0.9cm] (task) {\shortstack{任务定义\\$X,Y,\ell$}};
 \node[draw, rounded corners, fill=green!12, right=of task, minimum width=2.35cm, minimum height=0.9cm] (model) {\shortstack{模型\\$f_\theta$}};
 \node[draw, rounded corners, fill=yellow!20, right=of model, minimum width=2.35cm, minimum height=0.9cm] (train) {\shortstack{训练\\$\argmin_\theta$}};
 \node[draw, rounded corners, fill=red!10, right=of train, minimum width=2.15cm, minimum height=0.9cm] (deploy) {\shortstack{推理\\与部署}};
 \draw[->, thick] (data)--(task); \draw[->, thick] (task)--(model); \draw[->, thick] (model)--(train); \draw[->, thick] (train)--(deploy);
 \draw[->, thick, dashed] (deploy.north) .. controls +(0,0.72) and +(0,0.72) .. node[above,font=\scriptsize]{反馈、漂移、错误分析} (data.north);
\end{tikzpicture}
\caption{工业数据智能的训练、部署与反馈：部署反馈会重新影响数据和任务定义。}
\label{fig:foundations-learning-loop}
\end{figure}

\section{从问题定义到模型计算}

本章确定了输入、标签、模型和训练目标之间的关系。经验风险描述训练数据上的平均损失，总体风险描述目标分布上的预测代价。归纳偏置和正则化能够约束学习，但其效果仍需要独立数据检验。

接下来需要把这些定义用于计算。第三章以线性回归说明怎样组织数据、求出参数，以及怎样用后验分布表达参数不确定性。第四章进一步比较不同任务的损失和评价指标，第五章讨论目标没有直接解或后验难以计算时的求解方法。
\section*{辅助学习材料}
\begingroup
\emergencystretch=3em
\begin{itemize}
\item \textbf{Linear Models}；作者/平台：scikit-learn Developers；英文；官方文档。重点阅读 Ordinary Least Squares、Ridge、Lasso 与 Logistic Regression 小节，配合本章的线性预测、正则化和分类损失理解参数约束与风险最小化。\href{https://scikit-learn.org/stable/modules/linear\_model.html}{在线阅读}。
\item \textbf{Cross-validation: evaluating estimator performance}；作者/平台：scikit-learn Developers；英文；官方文档。重点阅读训练/验证/测试划分、交叉验证和数据泄漏说明，帮助把本章的经验风险、部署风险与可复现实验边界落实为流程。\href{https://scikit-learn.org/stable/modules/cross\_validation.html}{在线阅读}。
\item \textbf{损失函数（Loss Function）}；知乎；中文解说。比较平方损失与交叉熵，核对本章预测对象；标题与主题据必应索引，原页受限。\href{https://zhuanlan.zhihu.com/p/261059231}{在线阅读}。
\item \textbf{损失函数（lossfunction）的全面介绍（简单易懂版）}；CSDN；中文博客。阅读 L1、MSE 与交叉熵说明，比较不同误差的代价；已核对公开页面标题与索引摘要。\href{https://blog.csdn.net/weixin\_57643648/article/details/122704657}{在线阅读}。
\item \textbf{2021 - 类神经网络训练不起来怎么办(四) 损失函数 (Loss) 也可能有影响}；李宏毅；bilibili；中文课程。选看第22集，对照本章损失与学习目标。2026年10月4日官方公开元数据显示整套课程累计播放约299.5万，并非本分集播放量；已核对目录与计数，未逐帧观看。\href{https://www.bilibili.com/video/BV1Wv411h7kN/?p=22}{观看分集}。
\end{itemize}
\endgroup

